% Generated draft: requires mathematical and source review.
\documentclass[11pt]{article}
\usepackage[margin=25mm]{geometry}
\usepackage{fontspec}
\setmainfont{Latin Modern Roman}
\usepackage{amsmath,amssymb,mathtools,unicode-math}
\setmathfont{Latin Modern Math}
\usepackage{enumitem,needspace,xurl,hyperref}
\setlength{\parindent}{0pt}
\setlength{\parskip}{5pt}
\setlength{\emergencystretch}{3em}
\newcommand{\N}{\mathbb{N}}
\newcommand{\Z}{\mathbb{Z}}
\newcommand{\Q}{\mathbb{Q}}
\newcommand{\R}{\mathbb{R}}
\newcommand{\C}{\mathbb{C}}
\newcommand{\F}{\mathbb{F}}
\newcommand{\A}{\mathbb{A}}
\newcommand{\PP}{\mathbb{P}}
\newcommand{\E}{\mathbb{E}}
\newcommand{\eps}{\varepsilon}
\newcommand{\dd}{\,\mathrm{d}}
\begin{document}
\section*{UnsolvedMath ID: 800314}
\textbf{Source problem code:} catalogue entry
\subsection{Definitions and mathematical statement}
Let \(\mathrm{S}^1=\R/\Z\), and let \(G=\mathrm{Diff}^{\omega}_+(\mathrm{S}^1)\) be the group, under composition, of orientation-preserving real-analytic diffeomorphisms of the circle. Write \(H^\bullet_{\mathrm{cont}}(G)\) for the continuous cohomology intended by the source, with its cup product, and denote its named Euler and Godbillon--Vey classes by \(e\) and \(\mathrm{GV}\), respectively.

The source asks whether these two classes generate the continuous cohomology as a graded algebra. Symbolically, writing \(k\) for the intended coefficient ring, the question is
\[
H^\bullet_{\mathrm{cont}}(G)\stackrel{?}{=}k[e,\mathrm{GV}].
\]
Here \(k[e,\mathrm{GV}]\) means the unital subalgebra generated by the two classes: every cohomology class should be expressible as a finite sum
\[
\sum_{(i,j)\in F} a_{ij}\,e^{\smile i}\smile\mathrm{GV}^{\smile j},
\qquad F\subseteq\Z_{\geq 0}^2\text{ finite},\quad a_{ij}\in k,
\]
where a zeroth cup power is the unit. This notation does not assert algebraic independence.

\textit{Unresolved conventions.} The supplied record does not specify a topology on \(G\), the coefficients and their action, the precise continuous-cohomology construction, or the degrees and definitions of the two named classes in that construction. The notation above retains these unspecified conventions; \(k\) is a placeholder, not a universally quantified coefficient ring. A fully specified mathematical statement requires these choices to be recovered from the source. The interpretation of ``generated'' as generation under cup product is also an interpretation of the supplied wording.
\subsection{Short English statement}
Navas asks whether every continuous cohomology class of the group of orientation-preserving real-analytic circle diffeomorphisms can be expressed using the Euler and Godbillon--Vey classes, coefficients, addition, and cup products. The supplied record leaves the topology and cohomological conventions unspecified and labels the question open; this draft does not independently verify that status.
\subsection{Sources}
\begin{itemize}
\item Local UnsolvedMath registry, supplied complete source record: permanent ID \(800314\), exact source problem code \texttt{catalogue entry}, title ``Group Actions on 1-Manifolds --- Question 31.'' The record attributes the question to Navas, with proposed year 2017, and supplies the status \texttt{open}.
\item Navas, \emph{Group Actions on 1-Manifolds: A List of Very Concrete Open Questions} (2017), source item recorded as ``Question 31 (\#30).'' Supplied source URL: \url{https://arxiv.org/abs/1712.06462}.
\item Supplied provenance records access on 2026-10-16 and extraction by \texttt{arxiv\_latex}. These are supplied provenance claims, not an independent source visit or verification. The local record was updated on 2026-10-16 and identifies upstream revision \texttt{d371ee917ac9647fbb06cd6fc529deec6f3dda82}.
\end{itemize}
\end{document}
